\clearpage
\SourceLeaf{299}
\GaussVSourceUnnumbered
\pdfbookmark[1]{Göttingische gelehrte Anzeigen. 1832 December 24.}{anzeigen-1832-december-24}
\begin{center}
\rule{.70\linewidth}{.4pt}\\[.7\baselineskip]
Göttingische gelehrte Anzeigen. 1832 December 24.\\[.7\baselineskip]
\rule{.70\linewidth}{.4pt}
\end{center}

In der Sitzung der Königlichen Societät am 15. December wurde von dem
Hrn. Hofr. \textsc{Gauss} eine Vorlesung gehalten:
\begin{center}
\textit{Intensitas vis magneticae terrestris ad mensuram absolutam revocata,}
\end{center}
\noindent
von deren Gegenstande hier ein Bericht zu geben ist.

Dass von den drei, die Äusserung des Erdmagnetismus an einem gegebenen Orte
bestimmenden Elementen, Declination, Inclination und Intensität, das erste am
frühesten, viel später das zweite, und das dritte erst in den neuesten Zeiten
Gegenstand der Beobachtungen und Forschungen geworden ist, erklärt sich
hauptsächlich aus dem Umstande, dass die Declination für Seefahrer und Geodäten
das unmittelbarste Interesse darbietet, und die Inclination ihr näher verwandt
geschienen haben mag, als die Intensität. Bei dem Naturforscher, als solchem,
ist das Interesse für alle drei Elemente ganz gleich: unsere Kenntniss von dem
Erdmagnetismus im Ganzen muss so lange unvollkommenes Stückwerk bleiben, als
nicht alle Zweige derselben mit gleicher Liebe gepflegt werden.

Die ersten Aufklärungen über die Intensität des Erdmagnetismus verdanken wir
Herrn \textsc{von Humboldt}, welcher auf allen seinem Reisen ein Hauptaugenmerk
darauf gerichtet, und eine sehr grosse Menge von Beobachtungen geliefert hat,
aus denen sich die allmähliche Abnahme dieser Intensität, von dem magnetischen
Aequator der Erde nach den magnetischen Polen zu, ergeben hat. Sehr
\clearpage
\SourceLeaf{300}
\GaussVSourceFolio{294}{L}{ANZEIGEN.}
\noindent
viele Beobachter sind seitdem in die Fussstapfen jenes grossen Naturforschers
getreten, und ein Schatz von Beobachtungen aus fast allen Theilen der
Erdoberfläche, wohin in neuester Zeit wissenschaftliche Reisende gekommen sind,
liegt vor, worauf der um die Kenntniss des Erdmagnetismus hochverdiente
\textsc{Hansteen} bereits den Versuch einer allgemeinen isodynamischen Karte
hat begründen können.

Die bei allen diesen Beobachtungen angewandte Methode besteht darin, dass
man an den Örtern, für welche man die Intensität des Erdmagnetismus unter sich
vergleichen will, eine und dieselbe Magnetnadel Schwingungen machen lässt,
und deren Dauer mit Schärfe abmisst. Diese Dauer ist zwar, bei sonst gleichen
Umständen, von der Grösse des Schwingungsbogens abhängig, jedoch so, dass sie,
wie klein auch der Bogen wird, nur einer bestimmten Grenze immer näher kommt,
die man schlechthin die Schwingungsdauer nennt, und auf welche die wirklich
beobachtete vermittelst der Kenntniss des Schwingungsbogens leicht reducirt
werden kann. Die Intensität des Erdmagnetismus ist so dem Quadrate der
Schwingungsdauer einer und derselben Nadel verkehrt, oder dem Quadrate der
Anzahl der Schwingungen in einer gegebenen Zeit direct proportional, und das
Resultat bezieht sich auf die ganze Kraft, oder auf den horizontalen Theil
derselben, je nachdem man die Nadel in der Ebene des magnetischen Meridians
um eine horizontale, oder in einer horizontalen Ebene um eine verticale Axe
hat schwingen lassen.

Offenbar ist die Zulässigkeit dieses Verfahrens gänzlich von der Voraussetzung
des unveränderten magnetischen Zustandes der gebrauchten Nadel abhängig.
Wenn eine zweckmässig magnetisirte und vorsichtig aufbewahrte Nadel aus gut
gehärtetem Stahl zu diesen Versuchen angewandt wird, und diese selbst keinen
zu langen Zeitraum umfassen, wird freilich die Gefahr einer bedeutenden
Veränderung im Zustande der Nadel selbst nicht sehr gross sein, und man
kann sich darüber um so mehr beruhigen, wenn man nach der Zurückkunft an
den ersten Ort daselbst dieselbe Schwingungsdauer wiederfindet: allein selbst
die Erfahrung lehrt, dass man auf einen solchen Erfolg nicht leicht rechnen
darf, und genau genommen, enthält selbst jene Beruhigung einen logischen
Zirkel. In der That ist längst bekannt, dass sowohl die Declination, als die
Inclination an einem bestimmten Orte keinesweges unveränderlich ist; beide
erleiden im Lauf der Zeit sehr grosse fortschreitende, so wie daneben nach
den Tages- und Jahreszeiten
\clearpage
\SourceLeaf{301}
\GaussVSourceFolio{295}{R}{INTENSITAS VIS MAGNETICAE TERRESTRIS AD MENSURAM ABSOLUTAM REVOCATA.}
\noindent
für feinere Beobachtungen sehr merkliche periodische Veränderungen; es lässt
sich daher nicht zweifeln, dass auch das dritte Element, die Intensität,
ähnlichen Änderungen unterworfen sein wird, ja, die periodischen Änderungen
in verschiedenen Tageszeiten lassen sich in feinern Beobachtungen bestimmt
nachweisen. Wenn man daher auch nach längerer Zeit an demselben Orte dieselbe
Schwingungszeit wiederfindet, so hat man doch durchaus keine Bürgschaft,
dass dies nicht einer zufälligen Compensation der Veränderungen, welche
die Intensität des Erdmagnetismus an diesem Orte und der magnetische Zustand
der Nadel selbst inzwischen erlitten haben, zuzuschreiben sei. Wenn man auch
zugibt, dass durch diesen Umstand die Sicherheit der comparativen Methode,
in sofern nur mässige Zwischenzeiten vorkommen, nur etwas vermindert, nicht
ganz aufgehoben wird, so ist doch klar, dass diese Methode alle Brauchbarkeit
verliert, wenn die Frage die ist, welche Veränderung die Intensität des
Erdmagnetismus an einem bestimmten Orte während eines sehr langen Zeitraums
erfahre, und dass diese doch in wissenschaftlicher Beziehung höchst
interessante Frage ganz unbeantwortet bleibt, wenn man nicht an die Stelle
jener bloss comparativen Methode eine andere setzen kann, welche die
Intensität des Erdmagnetismus auf ganz bestimmte, für sich feststehende,
jederzeit mit grösster Schärfe wieder nachzuweisende und von der Individualität
der angewandten Nadeln ganz unabhängige Einheiten zurückführt.

Es ist nicht schwer, die theoretischen Grundsätze, auf welchen eine solche
selbstständige Methode beruhen muss, anzugeben. Die Schwingungsdauer einer
bestimmten Nadel hängt von drei Grössen ab, der Intensität des Erdmagnetismus,
dem statischen Moment des freien Magnetismus in der Nadel und dem Moment
der Trägheit der Nadel: letzteres kann leicht durch schickliche Methoden
ausgemittelt werden, und so ergibt sich aus der beobachteten Schwingungsdauer
nicht die Grösse der Intensität des Erdmagnetismus, sondern das Product dieser
Grösse in das statische Moment des freien Magnetismus der Nadel. Allein es
ist unmöglich diese beiden Factoren von einander zu trennen, wenn nicht
Beobachtungen einer ganz andern Art hinzukommen, die eine verschiedene
Combination derselben involviren: diesen Zweck erreichen wir, wenn wir eine
zweite Nadel zuziehen und dieselbe sowohl der Einwirkung des Erdmagnetismus
als der der ersten Nadel unterwerfen, um das Verhältniss dieser Kräfte
ausfindig zu machen. Diese beiden Wirkungen hängen zwar mit von dem
magnetischen Zustande der
\clearpage
\SourceLeaf{302}
\GaussVSourceFolio{296}{L}{ANZEIGEN.}
\noindent
zweiten Nadel ab, allein eine schickliche Einrichtung der Versuche verschafft
uns die Möglichkeit, diesen zu eliminiren, indem das \textit{Verhältniss}
beider Wirkungen desto mehr davon abhängig wird, je grösser die Entfernung
der beiden Nadeln von einander angenommen wird. Offenbar wird aber dabei
zugleich die Lage der magnetischen Axen der beiden Nadeln und der ihre
Mittelpunkte verbindenden geraden Linien gegen den magnetischen Meridian,
und der magnetische Zustand der ersten Nadel zu berücksichtigen sein,
und alles dies wird dem Calcul nicht unterworfen werden können, ohne
das Gesetz der Kraft zu kennen, welches zwei Elemente freien Magnetismus
auf einander ausüben, d. i. womit sie, je nachdem sie gleichnamig oder
ungleichnamig sind, einander abstossen oder anziehen. Schon
\textsc{Tobias Mayer} hatte die Vermuthung aufgestellt, dass dieses Gesetz
dasselbe sei, wie das der allgemeinen Gravitation, d. i. dass jene Kraft
im umgekehrten Verhältniss des Quadrats der Entfernung stehe; \textsc{Coulomb}
und \textsc{Hansteen} haben diese Vermuthung durch Versuche zu bestätigen
gesucht: durch die in vorliegender Abhandlung enthaltenen Versuche ist
sie ausser allen Zweifel gesetzt. Dieses Gesetz bezieht sich aber nur
auf die Elementarwirkung; die Berechnung der Totalwirkung eines
magnetischen Körpers auf einen andern wird zu einer rein mathematischen
Aufgabe, so bald die Art der Vertheilung des freien Magnetismus in diesen
Körpern vollständig bekannt ist, und bleibt daher von deren zufälliger
Individualität abhängig; allein je grösser der Abstand wird, desto
geringer wird der Einfluss dieser Individualität, und bei sehr grossen
Entfernungen kann man (wie eben aus jenem Grundgesetz von selbst folgt)
die Gesammtwirkung unter sonst gleichen Umständen dem Cubus der Entfernung
umgekehrt proportional setzen. Das Product dieses Cubus in den Bruch,
welcher das Verhältniss der Wirkung der ersten Nadel und der Wirkung des
Erdmagnetismus, auf die zweite Nadel, ausdrückt, wird sich daher bei immer
wachsenden Entfernungen einer bestimmten Grenze nähern; eine zweckmässige
Combination von Beobachtungen in verschiedenen schicklich gewählten
Entfernungen wird, mathematisch behandelt, die Grenze kennen lehren,
aus welcher das \textit{Verhältniss} derjenigen beiden Grössen sich
herleiten lässt, deren Product aus den beobachteten Schwingungszeiten
abgeleitet war: die Verbindung beider Resultate gibt dann offenbar
diese beiden Grössen selbst.

Die Versuche zur Vergleichung der Wirkungen des Erdmagnetismus und der
ersten Nadel, auf die an einem Faden aufzuhängende zweite, können auf
\clearpage
\SourceLeaf{303}
\GaussVSourceFolio{297}{R}{INTENSITAS VIS MAGNETICAE TERRESTRIS AD MENSURAM ABSOLUTAM REVOCATA.}
\GaussVSourceGathering{38}
\noindent
zwiefache Art eingerichtet werden, indem letztere entweder im Zustande der
Bewegung oder der Ruhe beobachtet werden kann. Das erstere geschieht am
vortheilhaftesten, indem man die erste Nadel in den magnetischen Meridian
der zweiten legt, wodurch die Dauer einer Schwingung der letztern entweder
grösser oder kleiner wird, je nachdem gleichnamige oder ungleichnamige
Pole einander zugekehrt sied: die Vergleichung der so veränderten
Schwingungsdauer mit der durch den blossen Erdmagnetismus bestimmten,
oder besser, die Vergleichung einer verlängerten mit einer verkürzten
(bei entgegengesetzten Lagen der ersten Nadel) führt dann leicht zu dem
gesuchten Verhältniss. Die zweite Art besteht darin, dass man die erste
Nadel so legt, dass ihre Einwirkung auf die zweite mit dem Erdmeridian
einen Winkel macht; der Ablenkungswinkel von dem Meridian, im Zustande
des Gleichgewichts, führt dann gleichfalls zur Kenntniss des verlangten
Verhältnisses, und auch hier ist es vortheilhafter, zwei entgegengesetzte
Ablenkungen, bei entgegengesetzten Lagen der ersten Nadel, unter sich zu
vergleichen. Die vortheilhafteste Lage dieser Nadel ist in einer durch
die Mitte der zweiten senkrecht auf den magnetischen Meridian gezogenen
geraden Linie. Übrigens kommt die erstere Art im wesentlichen mit
derjenigen überein, welche vor einigen Jahren von \textsc{Poisson}
vorgeschlagen ist; allein die bisher bekannt gewordenen Versuche einiger
Physiker, sie zur Anwendung zu bringen, sind entweder ganz missglückt,
oder können höchstens wie unvollkommene Annäherungen betrachtet werden.

Der Verfasser hat beide Arten vielfältig angewandt, und gefunden, dass
aus mehreren Gründen die zweite der ersten bei weitem vorzuziehen ist.

Die eigentliche Schwierigkeit liegt darin, dass in die beobachteten
Einwirkungen sich ausser dem Grenzwerthe noch andere Theile einmischen,
die von der Individualität der Nadeln abhängen. Jene Wirkung wird durch
eine Reihe dargestellt, die nach den negativen Potenzen des Abstandes
fortläuft, von der dritten anfangend, wo aber die folgenden Glieder sich
desto merklicher machen, je kleiner der Abstand ist. Man soll also aus
mehreren Beobachtungen diese folgenden Glieder eliminiren; allein bei
einiger Bekanntschaft mit der Eliminationstheorie überzeugt man sich
leicht, dass die unvermeidlichen Beobachtungsfehler der Zuverlässigkeit
des Resultats desto gefährlicher werden, je mehr Coëfficienten zu
eliminiren sind, so dass die Zahl derselben nur sehr mässig zu sein
braucht, um, aus jenem Grunde, die Rechnungsresultate gänzlich unbrauchbar
\clearpage
\SourceLeaf{304}
\GaussVSourceFolio{298}{L}{ANZEIGEN.}
\noindent
zu machen. Man hat daher keine Genauigkeit in den Resultaten zu erwarten,
wenn man nicht so grosse Distanzen anwendet, dass die Reihe sehr schnell
convergirt, und ein paar Glieder derselben zureichen. Allein dann sind
wieder die Wirkungen selbst sehr klein, also durch die bisherigen
Beobachtungsmittel nicht mit Schärfe zu bestimmen, und so erklärt sich
leicht das Misslingen der bisher angestellten Versuche.

So leicht sich also auch die Methoden, die Intensität des Erdmagnetismus
auf absolute Einheiten zurückzuführen, in der Theorie darstellen, so
misslich ist ihre Anwendung, so lange nicht den magnetischen Beobachtungen
eine viel grössere Schärfe verschafft wird, als sie bisher besassen.
Der Verfasser ist dadurch veranlasst, mehrere auf die Vervollkommnung
der Beobachtungsmittel abzweckende, schon vor vielen Jahren gefasste
Ideen zur Ausführung zu bringen, in der sichern Erwartung, dass die
magnetischen Beobachtungen zu einer beinahe, wo nicht ganz, eben so
grossen Schärfe zu bringen sind, wie die feinsten astronomischen.
Der Erfolg hat diese Erwartung nicht getäuscht, und zwei in der
Sternwarte aufgestellte Apparate, welche zu den zum Theil in vorliegender
Abhandlung aufgeführten Versuchen gedient haben, lassen nichts zu
wünschen übrig, als ein angemessenes gegen die Einwirkung von nahem
Eisen und Luftzug völlig geschütztes Local.

Es ist hier nicht der Ort zu einer vollständigen Beschreibung dieser
Apparate und ihrer Leistungen: wir glauben jedoch, den Naturforschern
durch Mittheilung der wesentlichsten Momente einen Dienst zu erweisen.

Die von dem Verfasser gewöhnlich gebrauchten Nadeln (wenn man prismatische
Stäbe von solcher Stärke noch Nadeln nennen darf) sind fast einen Fuss
lang, und haben ein Gewicht von beinahe einem Pfund. Die Aufhängung
geschieht an einem $2\tfrac12$ Fuss langen ungedreheten Seidenfaden, der
aus 32 einfachen zusammengesetzt, selbst das doppelte Gewicht noch
sicher trägt; das obere Ende des Fadens ist drehbar, und die Drehung
wird an einem eingetheilten Kreise gemessen. Die Nadel trägt an ihrem
südlichen oder nördlichen Ende (je nachdem die Localität das eine oder
das andere bequemer macht) einen Planspiegel, dessen Ebene gegen die
magnetische Axe der Nadel durch zwei Correctionsschrauben, so genau
wie man will, senkrecht gestellt werden kann, obwohl unnöthig ist,
hierauf eine sehr ängstliche Sorgfalt zu wenden, da man, was daran
fehlt, durch die Beobachtungen selbst auf das schärfste messen,
und als Collimations-
\clearpage
\SourceLeaf{305}
\GaussVSourceFolio{299}{R}{INTENSITAS VIS MAGNETICAE TERRESTRIS AD MENSURAM ABSOLUTAM REVOCATA.}
\GaussVSourceGathering{38*}
\noindent
fehler in Rechnung bringen kann.
Die so frei schwebende Nadel befindet sich in einem hölzernen
cylindrischen Kasten, welcher ausser der kleinen Öffnung im Deckel,
durch welche der Faden geht, noch eine grössere an der Seite hat,
welche nur wenig höher und breiter ist, als der erwähnte Spiegel.

Dem Spiegel gegenüber ist ein Theodolith aufgestellt; die verticale Axe
desselben und der Aufhängungsfaden sind in demselben magnetischen Meridian,
und etwa 16 Pariser Fuss von einander entfernt. Die optische Axe des
Fernrohrs am Theodolith ist etwas höher als die Nadel, und in der
Verticalebene des magnetischen Meridians so abwärts geneigt, dass sie
gegen die Mitte des Spiegels an der Nadel gerichtet ist.

An dem Stativ des Theodolithen ist eine 4 Fuss lange in einzelne Millimeter
getheilte horizontale Scale befestigt, die mit dem magnetischen Meridian
einen rechten Winkel macht; derjenige Punkt der Skale, welcher mit der
optischen Axe des Fernrohrs in Einer Verticalebene liegt, und hier Kürze
wegen der Nullpunkt heissen mag, wird durch einen von der Mitte des
Objectivs herabhängenden, mit einem Gewicht beschwerten feinen Goldfaden
bezeichnet; die Skale ist in einer solchen Höhe, dass das Bild eines
Theils derselben im Spiegel durch das Fernrohr erscheint, dessen Ocular
zum deutlichen Sehen auf die Entfernung dieses Bildes gestellt ist.
Auf der entgegengesetzten Seite von der Nadel ist in derselben
Verticalebene und in einer Entfernung vom Fernrohre, welche der jenes
Bildes gleich ist, eine Marke befestigt, welche dazu dient, jeden
Augenblick die unverrückte Stellung des Theodolithen prüfen zu können.

Es erhellt nun leicht, dass wenn obige Bedingungen genau erfüllt sind,
das Bild des Nullpunktes der Skale genau auf der optischen Axe des
Fernrohrs erscheinen muss, und dass, in sofern an dem Platze des
Theodolithen ein Gegenstand von bekanntem Azimuth sichtbar ist, man
mit Hülfe dieses Instruments sogleich die absolute magnetische
Declination erhalten kann. Fehlt dagegen an jenen Bedingungen etwas,
so wird, allgemein zu reden, nicht das Bild des Nullpunkts, sondern
das eines andern Punktes der Skale auf der optischen Axe erscheinen,
und wenn die horizontale Entfernung der Skale vom Spiegel genau
gemessen ist, wird der Betrag der Skalentheile leicht auf den
entsprechenden Winkel reducirt, und jenes erhaltene Resultat
corrigirt werden können. Der Betrag des Collimationsfehlers des
Spiegels kann mit grösster Schärfe und Leichtigkeit durch Umlegen
der Nadel in ihrem Träger (dass die obere Seite zur untern
\clearpage
\SourceLeaf{306}
\GaussVSourceFolio{300}{L}{ANZEIGEN.}
\noindent
wird) ausgemittelt werden. Bei den aufgestellten Apparaten beträgt
Ein Skalentheil nahe 22 Secunden, und ein nur etwas geübtes Auge
theilt ein solches Intervall noch leicht in 10 Theile.

Mit diesen Vorrichtungen bestimmt man also die Richtung der Nadel
und ihre Veränderungen auf das schärfste. Man hat gar nicht nöthig,
stets zu warten, bis sie zur Ruhe gekommen ist, da die beiden
Elongationen rechts und links sich mit äusserster Schärfe beobachten
lassen, und ihre Combination, gehörig behandelt, den entsprechenden
Ruhepunkt mit derselben Schärfe gibt. In den Vormittagsstunden, wo
die tägliche Variation am schnellsten ist, kann man diese beinahe
von Einer Zeitminute zur andern verfolgen.

Nicht minder gross ist der Gewinn dieser Einrichtung für die
Beobachtung der Schwingungsdauer. Das Vorübergehen des Verticalfadens
im Fernrohr vor einem bestimmten Punkte der Skale (eigentlich ists
umgekehrt), lässt sich, selbst wenn die ganze Ausweichung nur wenige
Minuten beträgt, mit einer solchen Schärfe beobachten, dass man bei
gehöriger Aufmerksamkeit niemals um ein ganzes Zehntheil einer
Zeitsecunde ungewiss bleibt. Die beträchtliche Dauer einer Schwingung
(bei den am kräftigsten magnetisirten Nadeln etwa 14 Secunden) und
die grosse Langsamkeit, mit welcher der Schwingungsbogen abnimmt,
gewähren hiebei noch andere höchst schätzbare Vortheile. Man braucht
nur ein paar Schwingungen beobachtet zu haben, um die Dauer Einer
Schwingung schon so scharf zu kennen, dass man dann die Nadel sich
selbst überlassen darf, und doch wenn man nach einer oder selbst
mehreren Stunden wieder hinzukommt, über die Anzahl der Schwingungen,
welche die Nadel in der Zwischenzeit gemacht hat, durchaus nicht
ungewiss ist. Man kann mit so kleinen Schwingungen anfangen (etwa
mit so grossen, wie die sind, bei denen man sonst aufzuhören pflegt),
dass die (übrigens äusserst leicht zu berechnende) Reduction auf
unendlich kleine Schwingungen fast unmerklich wird, und doch sind
dann nach 6 und mehreren Stunden die Schwingungen noch immer gross
genug, um die Antritte mit aller nöthigen Schärfe beobachten zu können.

Zeigen sich in den Beobachtungen zuweilen noch Anomalien, welche
aber stets so klein sind, dass sie bei den frühern Einrichtungen
gar nicht erkennbar gewesen sein würden, so sind solche einzig
dem in dem jetzigen Locale nicht immer \textit{ganz} zu vermeidenden
Luftzuge zuzuschreiben. Sie würden fast ganz wegfallen, wenn die
Öffnung des Kastens mit einem Planglase verschlossen würde,
\clearpage
\SourceLeaf{307}
\GaussVSourceFolio{301}{R}{INTENSITAS VIS MAGNETICAE TERRESTRIS AD MENSURAM ABSOLUTAM REVOCATA.}
\noindent
welches aber eine sehr grosse Vollkommenheit haben müsste. Dem
Verfasser stand bisher ein solches nicht zu Gebote, und jedenfalls
würde damit immer ein unangenehmer Lichtverlust verbunden sein.

Zu den bisher bemerkten Vortheilen dieser Einrichtungen kann man
noch den hinzufügen, dass der Beobachter stets in einer grossen
Entfernung von der Nadel bleibt, während er derselben bei den
früheren Verfahrungsarten sehr nahe kommen muss, und so, auch wenn
sie ganz in einem Glaskasten eingeschlossen ist, durch seine eigene
Wärme, durch die Wärme einer Beleuchtungslampe, oder durch Eisen
oder selbst Messing, welches er vielleicht bei sich führt, auf
die Nadel störend einzuwirken Gefahr läuft.

Der Vortheil, welchen starke schwere Nadeln, deren sich der
Verfasser ausschliesslich bedient, darbieten, ist so einleuchtend,
dass man es unbegreiflich finden muss, dass man sich zu den meisten
magnetischen Beobachtungen, namentlich für die Schwingungsdauer,
bisher immer nur äusserst kleiner Nadeln bedient hat. Es würde
vielmehr vortheilhaft sein, die von dem Verfasser bisher angewandten
Dimensionen noch weit zu überschreiten, was auch schon eine
versuchsweise gebrauchte Nadel von mehr als zwei Pfund Gewicht
bestätigt hat. Der Verfasser ist überzeugt, dass bei Anwendung
von vier- oder sechspfündigen Nadeln, wobei kleine Luftbewegungen
keinen merklichen Einfluss mehr haben werden, die magnetischen
Beobachtungen eine Schärfe erhalten können, die der der feinsten
astronomischen Beobachtungen durchaus nicht nachsteht. Freilich
muss man dann noch viel stärkere Aufhängungsfäden anwenden, deren
Torsion eine grössere Reaction ausüben wird; allein dies ist ganz
und gar kein Grund dagegen, da, für feine Resultate, die Torsionskraft
des Fadens doch nie ignorirt werden darf, sondern vielmehr, was
auch gar keine Schwierigkeit hat, jederzeit genau mit in Rechnung
gebracht werden muss.

Die beschriebenen Apparate dienen ausser dem Hauptzweck noch zu
einem andern, der, obgleich er mit jenem nicht in unmittelbarer
Verbindung steht, hier doch mit einigen Worten erwähnt werden
mag. Sie sind nemlich die schärfsten und bequemsten Galvanometer,
sowohl für die stärksten als für die schwächsten Kräfte eines
galvanischen Stroms, und es wird gar keine Schwierigkeit haben,
auch diese Messungen auf \textit{absolute} Maasse zurückzuführen.
Um die stärksten Kräfte zu messen, braucht man nur den Leitungsdraht
in beträchtlicher Entfernung (wenigstens mehrere Fuss) unterhalb
oder oberhalb der
\clearpage
\SourceLeaf{308}
\GaussVSourceFolio{302}{L}{ANZEIGEN.}
\noindent
Nadel im magnetischen Meridian einfach hinzuführen; für sehr schwache
Kräfte verbindet man damit einen Multiplicator, welcher um den die
Nadel enthaltenden Kasten gewunden ist. Der Verfasser hat einige
Versuche mit einem Multiplicator von 68 Drahtwindungen, die eine
Drahtlänge von 300 Fuss geben, gemacht: hier bedarf es keiner grossen
Plattenpaare; ein Paar kleine Knöpfe, ja selbst die blossen Enden
von Drähten aus verschiedenem Metall in gesäuertes Wasser eingetaucht,
bringen einen Strom hervor, der sich in einer Bewegung des
Skalenbildes von vielen hundert Skalentheilen sichtbar macht;
bei Anwendung von ein Paar Platten von sehr mässiger Grösse
fliegt hingegen im Augenblick der Schliessung der Kette das ganze
Skalenbild pfeilschnell durch das Gesichtsfeld des Fernrohrs.
Man übersieht leicht, wie sich durch diese Mittel die Abmessungen
an galvanischen Strömungen mit einer Schärfe und Bequemlichkeit
machen lassen, wovon die bisherigen mühsamen Methoden vermittelst
beobachteter Schwingungszeiten weit entfernt bleiben; man kann hier,
mit buchstäblicher Wahrheit, die allmähliche und bekanntlich anfangs
schnelle Abnahme der Stärke eines Stroms von Secunde zu Secunde
verfolgen. Will man noch anstatt einer einfachen Nadel eine doppelte
(astatische) anwenden, so wird keine electro-magnetische Kraft zu
klein sein, um nicht noch mit äusserster Schärfe gemessen werden zu
können. Es eröffnet sich demnach hier dem Naturforscher ein weites
reiches Feld für die interessantesten Untersuchungen.

Was nun den eigentlichen Hauptinhalt der Abhandlung betrifft, nemlich
die Entwickelung der mathematischen Theorie; verschiedene dem
Verfasser eigenthümliche Verfahrungsarten, z. B. zur Ausmittelung
des Moments der Trägheit der schwingenden Nadel, unabhängig von
der Voraussetzung einer regelmässigen Gestalt; die zur Constatirung
des oben erwähnten Grundgesetzes für die magnetischen Wirkungen
angestellten Versuche; endlich die Versuche zur Bestimmung des
Werths der Intensität des Erdmagnetismus an hiesigem Orte: so
müssen wir deshalb auf die Abhandlung selbst verweisen. Nur
die letzten Resultate wollen wir hier noch summarisch anführen.

Schon vor der Einrichtung der beschriebenen Apparate hatte der
Verfasser eine grosse Menge von Versuchen an Nadeln von den
verschiedensten Dimensionen, bis zu dem Gewichte von einem
halben Loth herab, angestellt, deren Resultate zwar sämmtlich
den spätern ziemlich nahe kommen, aber, da sie auf viel
unvollkommnern Hülfsmitteln beruhen, und weil es überhaupt
unmöglich ist,
\clearpage
\SourceLeaf{309}
\GaussVSourceFolio{303}{R}{INTENSITAS VIS MAGNETICAE TERRESTRIS AD MENSURAM ABSOLUTAM REVOCATA.}
\noindent
mit kleinern Nadeln eine grosse Schärfe zu erhalten, nicht mehr
verdienen aufbewahrt zu werden. Dagegen mögen sämmtliche mit
Hülfe der beschriebenen Apparate bisher erhaltene Resultate
für die Intensität des horizontalen Theils des Erdmagnetismus
hier Platz finden.
\begin{center}
\begin{tabular}{c|l|l}
I    & Mai 21       & 1.7820\\
II   & Mai 24       & 1.7694\\
III  & Jun. 4       & 1.7713\\
IV   & Jun. 24---28 & 1.7625\\
V    & Jul. 23, 24  & 1.7826\\
VI   & Jul. 25, 26  & 1.7845\\
VII  & Sept. 9      & 1.7764\\
VIII & Sept. 18     & 1.7821\\
IX   & Sept. 27     & 1.7965\\
X    & Octob. 15    & 1.7860
\end{tabular}
\end{center}

Als Einheiten liegen hiebei das Millimeter, das Milligramm und die
Zeitsecunde zum Grunde: \textit{wie} aber das Maass jener Intensität
durch diese Einheiten bestimmt ist, kann hier nicht entwickelt
werden: übrigens bleiben die Zahlen dieselben, wenn die Raumeinheit
und die Gewichtseinheit (eigentlich Masseneinheit) in gleichem
Verhältnisse geändert werden. Diese Versuche unterscheiden sich
theils durch die dabei beobachtete geringere oder grössere Sorgfalt,
theils durch die angewandten Nadeln, theils durch die Plätze,
auf welche sie sich beziehen.

Die Versuche VII. VIII. IX sind in jeder Beziehung so sorgfältig
ausgeführt, wie es nur der Apparat in seiner jetzigen Gestalt
verstattet, namentlich sind auch die dabei vorkommenden Distanzen
mit mikroskopischer Schärfe gemessen. Bei den Versuchen IV. V. VI. X
sind einige Operationen mit etwas geringerer Sorgfalt ausgeführt,
und die drei ersten Versuche stehen in dieser Beziehung noch
weiter zurück.

Zu den acht ersten Versuchen haben zwar verschiedene, aber an
Grösse und Gewicht nicht sehr ungleiche Nadeln (Gewicht zwischen
400 und 440 Grammen) gedient; die Hauptnadel im Versuch X wiegt
1062 Gramme; der Versuch IX hingegen ist mit einer sehr viel
kleineren Nadel (Gewicht 55 Gramme) angestellt, bloss um zu sehen,
welche Genauigkeit, bei Beobachtung jeder sonstigen
\clearpage
\SourceLeaf{310}
\GaussVSourceFolio{304}{L}{ANZEIGEN.}
\noindent
Vorsichtsmassregel, sich mit einer so kleinen Nadel erreichen
lasse: die Zuverlässigkeit des Resultats aus diesem Versuche
ist demnach den übrigen weit nachzusetzen.

Die Versuche VII\kern.1em.\kern.1em.\kern.1em.\kern.1em X sind
an Einem und demselben Platze in der Sternwarte angestellt,
die frühern hingegen an andern Plätzen, theils in der Sternwarte,
theils in den Wohnzimmern des Verfassers. Von allen diesen
Versuchen konnte demnach kein eigentlich reines Resultat erhalten
werden, da das in Gebäuden, und vorzüglich in der Sternwarte selbst,
vorhandene Eisen, durch den Erdmagnetismus selbst magnetisch
geworden, auf die Nadel reagirt, und seinen Einfluss mit dem
des Erdmagnetismus vermischt. Die Plätze sind übrigens immer
so gewählt, dass weder feste noch bewegliche Eisenmassen in
der \textit{Nähe} waren: allein einflusslos sind auch die
entferntern gewiss nicht geblieben. Indessen darf man doch
nach der blossen Ansicht der verschiedenen Resultate vermuthen,
dass die aus fremdem Einfluss herrührende Modification des
Erdmagnetismus an keinem dieser Plätze den hundertsten Theil
des Ganzen übersteigt. Ein eigentlich reines, und der Schärfe,
welche die Methode an sich verträgt, angemessenes Resultat,
wird man aber nur in einem eigenen Local, wo alles Eisen
gänzlich entfernt ist, erhalten können.

Um die Intensität der \textit{ganzen} Kraft des Erdmagnetismus
zu erhalten, müssen die gefundenen Zahlen noch mit der
Secante der Inclination multiplicirt werden. Der Verfasser
beabsichtigt, auch dieses Element in Zukunft nach eigenthümlichen
Methoden zu behandeln: einstweilen hat er am 23. Junius mit
dem Inclinatorium des physicalischen Cabinets $68^0\,22'\,52''$
gefunden, welches Resultat indessen, da die Beobachtung in
der Sternwarte, also nicht frei von fremdem Einfluss,
angestellt ist, leicht um mehrere Minuten dadurch
verändert sein kann.
\begin{center}
\rule{.25\linewidth}{.4pt}
\end{center}
